% Tabelas elaboradas a partir dos dados experimentais coletados.
% Os níveis de RSSI são apresentados na unidade bruta informada pelo módulo.

\newcommand{\TabelaAlcanceDistancia}{%
\begin{table}[htbp]
    \centering
    \caption{Síntese descritiva das medições por distância horizontal.}
    \label{tab:resultado-alcance-distancia}
    \small
    \setlength{\tabcolsep}{7pt}
    \begin{tabular}{cccc}
        \toprule
        \textbf{\shortstack{Distância\\(m)}} &
        \textbf{\shortstack{Detecções\\($n/N$)}} &
        \textbf{\shortstack{Taxa de\\leitura}} &
        \textbf{IC 95\% de Wilson} \\
        \midrule
        0,5 & 15/15 & 100,0\% & 79,6\% a 100,0\% \\
        1,0 & 85/105 & 81,0\% & 72,4\% a 87,3\% \\
        1,5 & 47/60 & 78,3\% & 66,4\% a 86,9\% \\
        2,0 & 67/105 & 63,8\% & 54,3\% a 72,4\% \\
        5,0 & 0/90 & 0,0\% & 0,0\% a 4,1\% \\
        \bottomrule
    \end{tabular}
    \fonte{Elaborado pelo autor a partir dos dados experimentais (2026).}
\end{table}%
}

\newcommand{\TabelaOrientacao}{%
\begin{table}[htbp]
    \centering
    \caption{Síntese descritiva das medições por orientação da etiqueta.}
    \label{tab:resultado-orientacao}
    \small
    \setlength{\tabcolsep}{7pt}
    \begin{tabular}{cccc}
        \toprule
        \textbf{\shortstack{Orientação\\($^\circ$)}} &
        \textbf{\shortstack{Detecções\\($n/N$)}} &
        \textbf{\shortstack{Taxa de\\leitura}} &
        \textbf{IC 95\% de Wilson} \\
        \midrule
        0 & 115/165 & 69,7\% & 62,3\% a 76,2\% \\
        45 & 60/105 & 57,1\% & 47,6\% a 66,2\% \\
        90 & 39/105 & 37,1\% & 28,5\% a 46,7\% \\
        \bottomrule
    \end{tabular}
    \fonte{Elaborado pelo autor a partir dos dados experimentais (2026).}
\end{table}%
}

\newcommand{\TabelaMaterial}{%
\begin{table}[htbp]
    \centering
    \caption{Resultado experimental por material de fixação.}
    \label{tab:resultado-material}
    \small
    \setlength{\tabcolsep}{4pt}
    \begin{tabular}{lcccc}
        \toprule
        \textbf{Material de fixação} &
        \textbf{\shortstack{Detecções\\($n/N$)}} &
        \textbf{\shortstack{Taxa de\\leitura}} &
        \textbf{IC 95\%} &
        \textbf{\shortstack{RSSI bruto\\média $\pm$ DP}} \\
        \midrule
        Papelão & 5/5 & 100,0\% & 56,6\% a 100,0\% & 183,0 $\pm$ 4,5 \\
        Metal direto & 0/5 & 0,0\% & 0,0\% a 43,4\% & Sem leitura \\
        Metal com espaçador & 5/5 & 100,0\% & 56,6\% a 100,0\% & 182,6 $\pm$ 4,2 \\
        \bottomrule
    \end{tabular}
    \fonte{Elaborado pelo autor a partir dos dados experimentais (2026).}
\end{table}%
}

\newcommand{\TabelaPorEtiqueta}{%
\begin{table}[htbp]
    \centering
    \caption{Desempenho por conjunto formado pela etiqueta e pelo suporte no subconjunto com registros individuais.}
    \label{tab:resultado-por-tag}
    \small
    \setlength{\tabcolsep}{4pt}
    \begin{tabular}{lrrrr}
        \toprule
        \textbf{\shortstack{Etiqueta\\e suporte}} &
        \textbf{\shortstack{Distância\\detecções ($n/N$)}} &
        \textbf{\shortstack{Taxa de\\leitura}} &
        \textbf{\shortstack{Orientação\\detecções ($n/N$)}} &
        \textbf{\shortstack{Taxa de\\leitura}} \\
        \midrule
        TAG1 (papelão) & 15/20 & 75,0\% & 13/15 & 86,7\% \\
        TAG2 (madeira) & 17/20 & 85,0\% & 13/15 & 86,7\% \\
        TAG3 (plástico) & 12/20 & 60,0\% & 11/15 & 73,3\% \\
        \bottomrule
    \end{tabular}
    \fonte{Elaborado pelo autor a partir dos dados experimentais (2026).}
\end{table}%
}

\newcommand{\TabelaInventarioRodadas}{%
\begin{table}[htbp]
    \centering
    \caption{Cobertura por rodada no inventário experimental em sala.}
    \label{tab:resultado-inventario-rodadas}
    \small
    \setlength{\tabcolsep}{6pt}
    \begin{tabular}{ccccc}
        \toprule
        \textbf{Rodada} &
        \textbf{\shortstack{Etiquetas\\detectadas}} &
        \textbf{Cobertura} &
        \textbf{\shortstack{Rodada\\completa}} &
        \textbf{\shortstack{Leituras\\externas}} \\
        \midrule
        1 & 4 & 80,0\% & Não & 0 \\
        2 & 4 & 80,0\% & Não & 0 \\
        3 & 4 & 80,0\% & Não & 0 \\
        4 & 4 & 80,0\% & Não & 0 \\
        5 & 3 & 60,0\% & Não & 0 \\
        \bottomrule
    \end{tabular}
    \fonte{Elaborado pelo autor a partir dos dados experimentais (2026).}
\end{table}%
}

\newcommand{\TabelaInventarioResumo}{%
\begin{table}[htbp]
    \centering
    \caption{Resumo do inventário experimental.}
    \label{tab:resultado-inventario}
    \small
    \begin{tabular}{p{6.2cm}r}
        \toprule
        \textbf{Métrica} & \textbf{Valor} \\
        \midrule
        Rodadas analisadas & 5 \\
        Etiquetas presentes por rodada & 5 \\
        Cobertura média & 76,0\% \\
        Desvio-padrão da cobertura & 8,9 pontos percentuais \\
        IC 95\% da cobertura média & 64,9\% a 87,1\% \\
        Faixa de cobertura & 60,0\% a 80,0\% \\
        Taxa de rodadas completas & 0/5 = 0,0\% \\
        IC 95\% da taxa de rodadas completas & 0,0\% a 43,4\% \\
        Leituras externas totais & 0 \\
        Leituras repetidas totais & 54 \\
        Leituras repetidas médias por rodada & 10,8 \\
        Leituras válidas médias por rodada & 14,6 \\
        \bottomrule
    \end{tabular}
    \fonte{Elaborado pelo autor a partir dos dados experimentais (2026).}
\end{table}%
}
